\addchap{\lsPrefaceTitle}
\markdouble{\lsPrefaceTitle}
 
%\chapter*{Preface}
%\addcontentsline{toc}{chapter}{Preface}
%\markboth{PREFACE}{PREFACE}

This book is a revised and shortened version of my PhD dissertation which I defended on June 28th 2024 at Goethe University Frankfurt. It aims at an analysis of the Japanese DP in line with the cartographic research program. On the one hand, the study presented here emphasizes the validity of the cartographic assumptions and their applicability to languages outside the Indo-European cosmos, but at the same time it shows that Japanese forces us to rethink many previously made claims and allows for refinements of these assumptions and a better understanding of the DP in cross-linguistic perspective.

The audience I would like to address are linguists interested in the DP and nominal modification as well as scholars in the field of Japanese studies who want to take a closer look at the different types of modification and the internal structure of modifiers in Japanese and to see how these language-specific phenomena can be transferred into a cross-linguistic framework. I strongly believe that both parties can equally benefit from reading this book.

The book consists of 6 chapters. Chapter 1 introduces the scope of this book and the research questions to be answered and gives an overview over the main hypotheses of this book. It also presents an overview over the Japanese nominal domain and a note on data collection. Chapter 2 is concerned with the theoretical background, the research program \textit{Cartography}. Chapter 3 addresses the main hypotheses defended in this book. It presents in detail the assumed DP structure and discusses the implications and necessary amendments to be made in line with the Japanese data, compared to earlier studies. The second part of this chapter is dedicated to the monomoraic elements co-occurring with non-verbal modifiers, the widely-distributed element \textsc{no}, as well as the elements \textsc{i} and \textsc{na}, which solely occur with the two Japanese adjective groups. I propose that those elements are instances of a syntactic head \textsc{mod} which licenses elements in attributive position.

Chapters 4 and 5 are dedicated to the direct and indirect domain respectively. In Chapter 4, I will outline each group of direct modifiers based on criteria for direct modification, both cross-linguistic and language-specific. Following is a discussion on the lack of ordering restrictions, otherwise a prerequisite for direct modification. Chapter 5 starts with an account of the complex elements co-occurring with Japanese modifiers. Then the size of indirect modifiers will be evaluated and a uniform CP status will be argued for. This chapter provides similar findings to the direct domain and closes with the assumption that more complex relative clauses are more prone to move to a higher position within the indirect domain.

The book closes with a summary and an outlook. Following are two appendices which contain information on the group of \textit{extended modifiers} and a detailed taxonomy of the group of \textit{no}-modifiers.

All Japanese examples throughout this book, if not quoted otherwise, were provided to me by native speakers of Japanese. Most examples were kindly provided to me with additional discussion by Ken Hiraiwa. Other native speakers who helped me with examples are: Jiro Inaba, Emiko Hayatsu, Chikako Takeuchi, Noboru Miyazaki.

